\section{Summary}

The experimental results on both the Oxford-IIIT Pet and \gls{TNBC}
(Crowd-Seg) datasets provide valuable insights into the performance
and limitations of our proposed approach. The model demonstrated
exceptional performance on the Oxford-IIIT Pet dataset, achieving
competitive Dice and Jaccard scores while successfully capturing
annotator reliability patterns. The model's ability to simultaneously
learn segmentation and reliability estimation proved particularly
valuable, providing insights into annotator behavior that traditional
approaches like STAPLE and MV cannot offer.

However, the experiments on the \gls{TNBC} dataset revealed significant
challenges in evaluating semantic segmentation models for
histopathological images when dealing with unreliable ground truth.
The morphological inconsistencies in the expert annotations, combined
with the extremely low label rates (1-12\%) across annotators, made
it difficult to objectively measure the model's performance. While
the model showed promising convergence in training, the validation
metrics remained unstable, suggesting that the current dataset may
not be suitable for rigorous evaluation of multi-annotator
segmentation approaches.

These findings highlight the need for more reliable ground truth
annotations in histopathological image segmentation datasets. Future
work should focus on developing datasets with more consistent and reliable
expert annotations, implementing stricter quality control measures for ground
truth generation, exploring alternative evaluation metrics that account for
annotation uncertainty, and investigating methods to handle extremely sparse
annotations more effectively. Recent work by \cite{LiEtAl2024} has addressed
the challenge of uncertainty quantification in medical image segmentation,
particularly focusing on inter-rater variability through the QUBIQ challenge.
Their comprehensive evaluation across various modalities (MRI, CT) and organs
(brain, prostate, kidney, pancreas) demonstrates the importance of ensemble
models and highlights the need for further research in 3D uncertainty
quantification methods, which could be particularly relevant for improving
the reliability of crowdsourced medical imaging annotations for segmentation
tasks.

Despite these challenges, the model's performance on the Oxford-IIIT
Pet dataset demonstrates its effectiveness in handling multiple
annotators with varying reliability levels. The model
successfully learned to:
\begin{itemize}
  \item Balance the contributions of different annotators based on
    their reliability
  \item Provide detailed spatial reliability information through
    pixel-wise estimation
  \item Maintain high segmentation accuracy even with missing labels
  \item Adapt to different noise levels in annotator inputs
\end{itemize}

These capabilities make the proposed approach particularly valuable
for real-world applications where multiple annotators with varying
expertise levels contribute to the dataset. The model's ability to
simultaneously learn segmentation and reliability estimation provides
a more comprehensive understanding of the annotation process, which
could be leveraged to improve the quality of future annotations and
dataset construction.